\section{Boundedness Analysis}\label{sec:analysis}

\begin{definition}[Expressions and Commands]\label{def:exp_com}
  \emph{Expressions} and \emph{commands} are defined by the following grammar,
  where $c$ is a literal constant of the language, $\mathit{aop}$ is an arithmetic
  operator and $\mathit{cop}$ is a comparison operator:
  \begin{align*}
    \exp::=&\ \mathit{c}\mid\mathit{v}\mid\field{\exp}{f}\mid\aop{\exp}{\exp}\mid\readmap{\exp}{\exp}{\exp}\\
    &\mid\new{C}{\exp^*}\mid\size{\exp}\mid\last{\exp}\mid\random{\exp}\\
    \com::=&\ \assign{\mathit{v}}{\exp}\mid\assign{\field{\exp}{f}}{\exp}\mid\writemap{\exp}{\exp}{\exp}
    \mid\require{\exp}{\exp}\\
    &\mid\push{\exp}{\exp}\mid\receive{\exp}{\exp}\mid\call{\exp}{m}{\exp^*}\\
    &\mid\ife{\exp}{\exp}{\com}{\com}\mid\for{\mathit{v}}{\exp}{\com}\mid\sequence{\com}{\com}
  \end{align*}
  Expressions and commands are silently annotated with the program point $p$ where they occur
  in the source code under analysis\footnote{For instance, a program point might be
  a line number and a character position inside that line,
  by its exact nature is irrelevant here.}. When useful, that annotation is made
  explicit as $\exp\at{p}$ or $\com\at{p}$.
\end{definition}

\begin{definition}[Program]\label{def:program}
  A \emph{program} is a collection of \emph{classes}.
  Each class $C$ has a superclass,
  some \emph{field signatures}, of the form $C.f$,
  and some \emph{methods}, of the form
  $C.m(p_1 t_1,\ldots,p_n t_n)\mathtt{\ \{}\com\mathtt{\}}$,
  where $\com$ is the \emph{body} of the method.
  A class that extends \<Contract>
  (possibly indirectly) is called a \emph{smart contract}.
  A \emph{constructor} is a method named like the class $C$
  that it belongs to.
\end{definition}
%
Fields and methods can have modifiers such as \<public>, \<private>,
\<final> and \<payable>, that we do not introduce formally and have the
usual meaning as, for instance, in Java or Solidity. For simplicity, we
assume that methods do not return any value.

We assume to analyze a statically strongly-typed language. Therefore, each
expression and command will be evaluated in a context (\emph{typing})
that specifies which variables are in scope and which is their type.
This typing is that inferred in the type-checking phase of the compiler.
%
\begin{definition}[Typing]\label{def:typing}
  A type is either \emph{primitive} (such as \<int>, \<boolean> \ldots)
  or a class name $C$. Generic class types (such as lists) are written as
  \<List<t\texttt{>}>, where \<t> is the type of the elements of the list.
  A \emph{typing} $\tau$ is a map from a finite set of variables
  to their static (\ieinitial declared) type. For any
  expression $\exp$, we write $\tau(\exp)$ for its \emph{type according to $\tau$},
  if its computation does not lead to a type error.
\end{definition}
%
The definition of $\tau(\exp)$ is not formalized here, since it is the usual recursive
definition of the type of an expression, used by any strongly-typed language.
We assume that each program point $p$ comes with an associated typing $\tau\at{p}$.
For any expression $\exp\at{p}$, we assume that $\tau\at{p}$ is \emph{correct} for $\exp$\ie the
computation of $\tau\at{p}(\exp)$ succeeds without type errors.
Such $\tau\at{p}$ is for instance that inferred by the compiler at $p$ during type-checking.

Given an expression or command in the program, it is possible to identify
which lists (and more generally,
which data structures) might be affected by the evaluation of the expression
or by the execution of the command.
Previous work has defined precise and complex static analyses for that
(see for instance the aliasing and side-effect analyses in~\cite{SpotoMP10}).
Here, for simplicity and lack of space,
we present a very simple and cheap side-effect analysis that, nevertheless,
is frequently precise enough on simple programs such as smart contracts.
%
\begin{definition}[Side-Effects Analysis]\label{def:side_effect}
  The set of element types of the lists potentially modified during the evaluation of an
  expression or the execution of a command in the program
  is overapproximated by the minimal fixpoint of the following recursive equations:
  %
  \begin{align*}
    \sideeffect{c}&=\sideeffect{v}=\varnothing\\
    \sideeffect{\new{C}{\exp_1,\ldots,\exp_n}}&=\sideeffect{\com}\cup\bigcup\limits_{i=1}^{n}\sideeffect{\exp_i}\\
    &\quad\text{where $\com$ is the body of the constructor}\\
    \sideeffect{\field{\exp}{f}}&=\sideeffect{\exp}\\
    \sideeffect{\aop{\exp_1}{\exp_2}}&=\sideeffect{\exp_1}\cup\sideeffect{\exp_2}\\
    &\quad\text{similarly, recursively, for the other expressions}
  \end{align*}
  \begin{align*}
    \sideeffect{\push{\exp_1}{\exp_2}\at{p}}&=\sideeffect{\exp_2}\cup\{t\}\quad\text{where $\tau\at{p}(\exp_1)=\<List<>t$\<\text{>}>}\\
    \sideeffect{\call{\exp_0}{m}{\exp_1,\ldots,\exp_n}\at{p}}&=\bigcup\limits_{j=1}^k\sideeffect{\com_j}\cup\bigcup\limits_{i=0}^{n}\sideeffect{\exp_i}\\
    &\quad\text{where $\com_1\cdots\com_k$ are the bodies of}\\
    &\quad\text{the methods possibly called at $p$}\\
    \sideeffect{\sequence{\com_1}{\com_2}}&=\sideeffect{\com_1}\cup\sideeffect{\com_2}\\
    \sideeffect{\assign{\mathit{v}}{\exp}}&=\sideeffect{\exp}\\
    &\quad\text{similarly, recursively, for the other commands.}
  \end{align*}
\end{definition}
%
We recall that a method call in object-oriented languages may actually call a subset
$\{M_1,\ldots,M_k\}$ of the redefinitions of its static target method, with bodies $\{\com_1,\ldots,\com_k\}$.
A sound overapproximation of $\{M_1,\ldots,M_k\}$ can always be inferred statically,
as the set of \emph{all} possible redefinitions. More precise analyses exist~\cite{TipP00},
but are probably overkill for minimalistic code such as smart contracts,
that very rarely redefine methods.

A \emph{path} is an expression rooted at a variable in scope, that follows
\<final> field dereferences or list accesses, or is an integer constant.
Paths have a type and a length. A constant $\kappa$ can be used to limit the maximal
length of a path. A relevant path has type list. While paths will be used
to define the definite aliasing analysis, relevant paths will be used in the
subsequent boundness analysis.
%
\begin{definition}[Paths]\label{def:path}
  Given $\kappa\in\nat$, the finite set $\mathbb{P}^\kappa\at{p}$ of the
  \emph{$\kappa$-paths at program point $p$} is the minimal set such that
  \begin{itemize}
  \item if $\kappa\ge 1$ and $v\in\dom(\tau\at{p})$ then $v\in\mathbb{P}^\kappa\at{p}$, it has type $(\tau\at{p})(v)$
    and $\length(v)=1$ (all variables in scope at $p$ are a $1$-path at $p$);
  \item if $\kappa\ge 1$ and $c\in\nat$ then $c\in\mathbb{P}^\kappa\at{p}$, it has type \<int>
    and $\length(c)=1$ (all integer constants are a $1$-path at $p$);
  \item if $\pi\in\mathbb{P}^\kappa\at{p}$ has type $C$ and $\length(\pi)<\kappa$,
    and if $\mathtt{f}$ is a \<final> field of type $t$ defined in class $C$,
    already initialized at $p$,
    then $\pi.f\in\mathbb{P}^\kappa\at{p}$, it has type $t$ and $\length(\pi.f)=\length(\pi)+1$
    (paths can be expanded by \<final> field dereference);
  \item if $\pi\in\mathbb{P}\at{p}$ has type \<List<$t$\texttt{>}> and $\length(\pi)<\kappa$
    then $\pi\<[any]>\in\mathbb{P}^\kappa\at{p}$, it has type $t$ and $\length(\pi\<[any]>)=\length(\pi)+1$
    (paths can be expanded by list access).
  \end{itemize}
  A $\kappa$-path is \emph{relevant} if its type is \<List<$t$\texttt{>}> for some $t$.
  The finite set of all relevant $\kappa$-paths at $p$ is written as $\mathbb{RP}^\kappa\at{p}$. 
\end{definition}
%
For instance, at line~\ref{ponzi1:require} in Fig.~\ref{fig:ponzi1},
$42$ is a $1$-path of type \<int>,
\<this> is a $1$-path of type \<Ponzi1>,
\<amount> is a $1$-path of type \<int>,
\<this.investors> is a $2$-path of type \<List<Contract\texttt{>}> and
\<this.investors[any]> is a $3$-path of type \<Contract>. Among them, only
\<this.investors> is relevant.

The $\kappa$ limit is essential to make our analysis finitely computable, as shown later.
Larger $\kappa$ lead to a more precise analysis,
in principle, but also potentially more expensive.
Therefore, $\kappa$ acts as a lever that selects the desired precision for the
analysis. For simplicity, in the examples that follow,
we assume $\kappa=3$ and just say \emph{relevant}, meaning $3$-relevant.

A path \emph{survives} the execution of some code\ie it keeps its value unchanged,
if it does not access any list that might be modified by that execution. Note how this
definition is simplified by the constraint that paths only dereference
\<final> fields (Def.~\ref{def:path}), which allows one to focus on side-effects on lists only.
%
\begin{definition}\label{def:survival}
  Given $\pi\in\mathbb{P}^\kappa\at{p}$ and an expression $\exp\at{p}$, then
  \emph{$\pi$ survives the evaluation of $\exp$} if, for all prefix $\pi'$ of $\pi$
  (\ieinitial  $\pi=\pi'\ldots$) of type \<List<$t'$\texttt{>}>
  and for all $t''\in\sideeffect{\exp}$, the types
  $t'$ and $t''$ have no common subtype.
  A similar definition states when \emph{$\pi$ survives the execution of a command $\com$}.
\end{definition}
%
For instance, the path $\pi=\<this.investors>$ at
line~\ref{ponzi1:expand} in Fig.~\ref{fig:ponzi1} does not survive the execution
of the command at that line, since $\pi'=\<this.investors>$
has type \<List<Contract\texttt{>}>,
$\{\<Contract>\}=\sideeffect{\<investors.push(caller)>}$
and \<Contract> has a common subtype with itself, obviously.
Instead, the path $\pi=\<last.sales>$ at line~\ref{lottery2:push_round}
survives the execution of the command at that line, since
the only prefix $\pi'$ of $\pi$ of type list is $\pi'=\pi$, of type
\<List<TicketBatchSale\texttt{>}>,
$\{\<Round>\}=\sideeffect{\<rounds.push(last)>}$ and
\<TicketBatchSale> and \<Round> have no common subtype.
All integer constants survive any expression and command.

\begin{definition}[Definite Aliasing Analysis]\label{def:definite_aliasing}
  An \emph{abstract definite aliasing state at $p$} is a map from each variable
  $v\in\dom(\tau\at{p})$
  to a set of paths that are definitely alias of $v$ (\ieinitial they have the same
  value as $v$). The \emph{abstract definite aliasing denotation} of an expression
  or command at program point $p$ is a map that transforms abstract definite aliasing
  pre-states at $p$ to abstract definite aliasing post-states at the program point after
  the expression or command, defined as:
  %
  \begin{align*}
    \da{c}(\sigma)&=\pair{\sigma}{\{c\}}\quad\text{if $\kappa\ge 1$}\\
    \da{v}(\sigma)&=\pair{\sigma}{\sigma(v)}\\
    \da{\field{\exp}{f}}(\sigma)&=\begin{cases}
    \pair{\sigma'}{\{\pi.f\mid\pi\in\val,\ \length(\pi)<\kappa\}} & \text{if $f$ is \<final>}\\
    \pair{\sigma'}{\varnothing} & \text{otherwise}
    \end{cases}\\
    &\quad\text{where $\da{\exp}(\sigma)=\pair{\sigma'}{\val}$}\\
    \da{\aop{\exp_1}{\exp_2}}(\sigma)&=\pair{\sigma_2}{\{\aop{c_1}{c_2}\mid c_1\in\val_1\cap\nat,\ c_2\in\val_2\cap\nat\}}\\
    &\quad\text{where $\da{\exp_1}(\sigma)=\pair{\sigma_1}{\val_1}$}\\
    &\quad\text{and $\da{\exp_2}(\sigma_1)=\pair{\sigma_2}{\val_2}$}
  \end{align*}
\end{definition}

The static analysis infers size change constraints about the containers
reachable from the state of the smart contracts, for each public API
method or constructor of the smart contract. It distinguishes between
pre-state and after-state by using primed variables for the after-state.
\textbf{This is a sort of simplified path-length. The simplifying
  assumptions must be clarified. The idea is that one does not need
  the supporting analyses such as non-cyclicity or sharing, since the
  domain of analysis is very specific and lists are very well encapsulated.}

For the example in Fig.~\ref{fig:ponzi1}, the inferred approximations are:
%
\begin{description}
\item[\<Ponzi1()>:] $\<investors>'=1$
\item[\<invest()>:] $\<investors>'=\<investors>+1$
\end{description}
%
For the esample in Fig.~\ref{fig:ponzi2} they are (\textbf{something about \<balances>? In any case, it is irrelevant for our proof}):
%
\begin{description}
\item[\<Ponzi2()>:] $\<investors>'=1$
\item[\<invest()>:] $\<investors>'=\<investors>+1$
\item[\<withdraw()>:] $\<investors>'=\<investors>$
\end{description}
%
For the example in Fig.~\ref{fig:ponzi3} they are:
%
\begin{description}
\item[\<Ponzi3()>:] $\<investors>'=1$
\item[\<invest()>:] $\<investors><100\wedge\<investors>'=\<investors>+1$
\end{description}
%
For the example in Fig.~\ref{fig:lottery1} they are:
%
\begin{description}
\item[\<Lottery1()>:] $\<rounds>'=1\wedge\<rounds[any].sales>'=0$
\item[\<buy()>:] $\<rounds>\le\<rounds>'\le\<rounds>+1\wedge\<rounds[any].sales>\le\<rounds[any].sales>'\le\<rounds[any].sales+1>$
\item[\<drawWinner()>:] $\<rounds>'=\<rounds>\wedge\<rounds[any].sales>'=\<rounds[any].sales>$
\end{description}
%
For the example in Fig.~\ref{fig:lottery2} they are (\textbf{automatically inferring the subsequent constraints might not be obvious}):
%
\begin{description}
\item[\<Lottery2()>:] $\<rounds>'=1\wedge\<rounds[any].sales>'=0$
\item[\<buy()>:] $\<rounds>\le\<rounds>'\le\<rounds>+1\wedge\<rounds[any].sales>\le\<rounds[any].sales>'\le\max(5000,\<rounds[any].sales>)$
\item[\<drawWinner()>:] $\<rounds>'=\<rounds>\wedge\<rounds[any].sales>'=\<rounds[any].sales>$
\end{description}

Such constraints are then used to build class invariant predicates,
stating that there exists an upper bound to the size of the containers
such that, if the containers are bounded before the call to the API,
they are bounded also after the call to the API.

For instance, for \<Ponzi1> the predicate is:
%
\[
  \exists k.\left(\begin{array}{c}
    (\mathit{true}\wedge\<investors>'=1)\Rightarrow\<investors>'\le k\\
    \wedge\\
    (\<investors>\le k\wedge\<investors'>=\<investors>+1)\Rightarrow\<investors>'\le k
  \end{array}\right)
\]
%
that a theorem prover should reject or fail to prove that it holds. Therefore,
the analyzer would issue a boundedness warning.

For \<Ponzi2>, the predicate is:
%
\[
  \exists k.\left(\begin{array}{c}
    (\mathit{true}\wedge\<investors>'=1)\Rightarrow\<investors>'\le k\\
    \wedge\\
    (\<investors>\le k\wedge\<investors>'=\<investors>+1)\Rightarrow\<investors>'\le k\\
    \wedge\\
    (\<investors>\le k\wedge\<investors>'=\<investors>)\Rightarrow\<investors>'\le k\\
  \end{array}\right)
\]
%
that a theorem prover should reject or fail to prove that it holds. Therefore,
the analyzer would issue a boundedness warning.

For \<Ponzi3>, the predicate is:
%
\[
  \exists k.\left(\begin{array}{c}
    (\mathit{true}\wedge\<investors>'=1)\Rightarrow\<investors>'\le k\\
    \wedge\\
    (\<investors>\le k\wedge\<investors><100\wedge\<investors>'=\<investors>+1)\\
    \Rightarrow\<investors>'\le k
  \end{array}\right)
\]
%
that a theorem prover should be able to prove true, for instance by taking
$k=100$ (or $k=200$\ldots). Therefore, the boundedness analyzer would not issue
any warning.

For \<Lottery1>, there is a predicate for \<rounds> and another for
\<rounds[any].sales>. The former is irrelevant, since \<rounds> is not
iterated. The latter is (by writing \<ras> for \<rounds[any].sales>):
%
\[
  \exists k.\left(\begin{array}{c}
    (\mathit{true}\wedge\<ras>'=0)\Rightarrow\<ras>'\le k\\
    \wedge\\
    (\<ras>\le k\wedge\<ras>\le\<ras>'\le\<ras>+1)\Rightarrow\<ras>'\le k\\
    \wedge\\
    (\<ras>\le k\wedge\<ras>'=\<ras>)\Rightarrow\<ras>'\le k
  \end{array}\right)
\]
%
that a theorem prover should reject or fail to prove that it holds. Therefore,
the analyzer would issue a boundedness warning.

For \<Lottery2>, the predicate for \<rounds[any].sales> is:
%
\[
  \exists k.\left(\begin{array}{c}
    (\mathit{true}\wedge\<ras>'=0)\Rightarrow\<ras>'\le k\\
    \wedge\\
    (\<ras>\le k\wedge\<ras>\le\<ras>'\le\max(5000,\<ras>))\Rightarrow\<ras>'\le k\\
    \wedge\\
    (\<ras>\le k\wedge\<ras>'=\<ras>)\Rightarrow\<ras>'\le k
  \end{array}\right)
\]
%
that a theorem prover should be able to prove true, for instance by taking
$k=5000$ (or $k=8000$\ldots). Therefore, the boundedness analyzer would not
issue any warning.

\textbf{Which theorem prover is able to deal with such constraints?}
